\documentclass[11pt,a4paper]{article}
\usepackage[utf8]{inputenc}
\usepackage[T1]{fontenc}
\usepackage[french]{babel}
\usepackage{lmodern}
\usepackage{geometry}
\geometry{a4paper, top=24mm, bottom=24mm, left=24mm, right=24mm}
\usepackage{xcolor}
\definecolor{NavyBlue}{HTML}{0A2540}
\definecolor{CobaltBlue}{HTML}{0055B3}
\definecolor{DarkSlate}{HTML}{2D3748}
\definecolor{LightGray}{HTML}{F7FAFC}
\usepackage{tcolorbox}
\usepackage{enumitem}
\usepackage{booktabs}
\usepackage{tabularx}
\usepackage{hyperref}
\hypersetup{colorlinks=true, linkcolor=CobaltBlue, urlcolor=CobaltBlue}

\title{\vspace{-1cm}\Huge\textbf{\textcolor{NavyBlue}{Mémoire PRO-ENT5 — Étape 3}}\\[6pt]
\Large\textbf{\textcolor{CobaltBlue}{Critique et Analyse Réflexive}}}
\author{\textbf{Jérémy SOLER} — Apprenti Architecte Systèmes d'Information\\
SAS Everbloo / Metis Digital — ETNA (Titre d'Expert en Ingénierie Informatique)}
\date{Année académique 2025 -- 2026}

\begin{document}
\maketitle

\vspace{0.2cm}
\begin{tcolorbox}[colback=LightGray, colframe=NavyBlue, arc=4pt, boxrule=1.1pt]
\centering\textbf{Problématique d'Ingénierie :}\\
\textit{« Comment garantir la cohérence transactionnelle et tarifaire d'un moteur de réservation aérien agrégeant des protocoles hybrides (GDS et NDC) ? »}
\end{tcolorbox}

\section*{Temps 1 — Analyse Critique et Confrontation des Approches}

\subsection*{1. Confrontation des Options d'Architecture}
Face aux contraintes du secteur aérien, plusieurs approches techniques ont été évaluées :
\begin{itemize}[itemsep=2pt]
    \item \textbf{REST/JSON Canonique vs GraphQL :} L'utilisation de GraphQL semblait séduisante pour limiter l'over-fetching sur les offres de vol volumineuses. Cependant, l'absence de mise en cache HTTP native et le fait qu'aucun transporteur aérien ne fournisse d'API GraphQL ont conduit à privilégier une architecture REST/JSON avec des DTOs typés TypeScript.
    \item \textbf{Flux Synchrone Orchestré vs Kafka Event-Driven :} Pour l'émission d'un billet d'avion, le client et le transporteur exigent une confirmation synchrone immédiate ($<$ 3\,s) pour garantir le prix. L'asynchronisme total via Kafka a donc été écarté pour le cycle direct de checkout, tout en restant pertinent pour le traitement différé des notifications de changement d'horaires (OCN).
    \item \textbf{Adaptateur Modulaire vs Couche d'Intégration Point-à-Point :} L'intégration ad-hoc aurait créé un code spaghetti ingérable. L'introduction du patron Adaptateur a garanti une étanchéité stricte entre les connecteurs externes et la logique métier.
\end{itemize}

\subsection*{2. Preuves de Terrain et Résultats Chiffrés}
L'architecture déployée a été éprouvée sur plus de 100 agences :
\begin{itemize}[itemsep=2pt]
    \item Taux d'échec de création d'ordre NDC réduit à $<$ 0,1\,\% (hors pannes serveur compagnies) ;
    \item Réduction de 28\,\% du temps de réponse moyen de shopping (de 4,2\,s à 3,0\,s) ;
    \item Élimination totale des litiges ADM pour double facturation de segments passifs.
\end{itemize}

\subsection*{3. Position Finale Défendue}
La position défendue repose sur une architecture orientée services (SOA) découplée, couplant un backend Node.js non-bloquant et des micro-frontends Nuxt 3 réactifs, sécurisée par une gestion de sessions OAuth2 avec file d'attente FIFO et un moteur de repricing interactif.

\vspace{0.4cm}
\section*{Temps 2 — Analyse Réflexive et Maturation Professionnelle}

\subsection*{1. Retour Critique sur l'Action et Gestion des Crises}
L'expérience m'a confronté aux réalités imprévues de la production : pannes d'API tierces un lundi matin, certificats TLS expirés sur les passerelles Sabre, blocages de transactions MySQL (deadlocks) et erreurs d'arrondis monétaires sur les taxes internationales.

Ces situations de stress m'ont appris à abandonner la panique pour une démarche méthodique : corrélation de logs Winston, diagnostic d'infrastructure et communication transparente auprès des agences partenaires.

\subsection*{2. Ce qui Serait Fait Différemment avec le Recul}
\begin{itemize}[itemsep=2pt]
    \item \textbf{Gestion des Versions Node.js :} Certains modules dépendent encore de Node 16. J'aurais dû imposer une migration complète vers Node 20 LTS plus tôt plutôt que de maintenir des scripts de compatibilité Pixi.
    \item \textbf{Monorepo et Sous-Modules Git :} L'utilisation de sous-modules Git a créé des frictions lors des déploiements concurrents. L'adoption d'un monorepo unifié avec Turborepo ou PNPM Workspaces aurait simplifié l'intégration continue.
\end{itemize}

\subsection*{3. Apports Professionnels et Personnels}
Cette immersion de 38 mois a transformé ma posture : je suis passé d'un rôle de développeur centré sur le code immédiat à celui d'un architecte capable de négocier avec le Product Management, de défendre des choix techniques face aux impératifs commerciaux, et d'assurer la pérennité d'un système critique.

\subsection*{4. Limites Actuelles et Pistes d'Approfondissement}
\begin{itemize}[itemsep=2pt]
    \item Intégration de la future norme \textbf{IATA ONE Order} pour unifier la commande et la billetterie ;
    \item Mise en place d'un traçage distribué complet avec \textbf{OpenTelemetry} pour superviser les latences des API transporteurs en temps réel.
\end{itemize}

\end{document}
